\documentclass[11pt]{article}
\usepackage[margin=1in]{geometry}
\usepackage{natbib}
\usepackage{booktabs}
\usepackage{hyperref}
\usepackage{amsmath}

\title{Vedaksha: A Clean-Room, Sub-Arcsecond Ephemeris Engine \\ for Computational Vedic Astrology}
\author{Amit Gupta \\ ArthIQ Labs LLC \\ \texttt{info@arthiq.net}}
\date{}

\begin{document}
\maketitle

\begin{abstract}
Vedaksha is a clean-room ephemeris and Vedic astrology (jyotish) computation engine written
in Rust. It computes planetary positions from the analytical VSOP87A and ELP/MPP02 series and
from the JPL DE440s SPK numerical kernel --- validated against an independent JPL DE441
ephemeris served by the NASA/JPL Horizons system --- and derives from them the constructs used in
Vedic astrology --- nakshatras, dashas, vargas (divisional charts), panchanga, shadbala, and
muhurta --- exposed through native Rust, Python, WebAssembly, and Model Context Protocol (MCP)
interfaces.
% TODO(Task 5): insert re-measured mean/max residual figure here once Validation section is written
Every algorithm is derived directly from a cited primary source rather than from an existing
implementation: the ephemeris and sidereal-conversion algorithms as shipped in the current
release contain no source code, header, or documentation from the GNU Affero General Public
Licensed (AGPL) Swiss Ephemeris library, on which prior Vedic astrology software has commonly
relied. (Three early sidereal-offset constants were briefly cross-referenced against a published
Swiss Ephemeris header solely to identify a numeric value; all three were independently re-derived
from cited primary sources in a subsequent clean-room pass, documented below.) This clean-room
provenance is not merely asserted: it is documented per data source, with fetch dates and,
where established, content hashes, in a public provenance ledger, and specific re-derivations
carry dated, auditable clean-room evidence trails. Vedaksha is released under the Business
Source License 1.1 (converting to Apache License
2.0 four years after each version's release) and is distributed via crates.io, PyPI, npm, and
GitHub.
\end{abstract}

\section{Introduction}

Software for computing Vedic astrology (jyotish) charts requires, as its foundation, a
planetary ephemeris of at least sub-arcsecond precision. In practice, the dominant ephemeris
used by existing Vedic astrology software --- commercial and open source alike --- is the Swiss
Ephemeris library, which is dual-licensed under the GNU Affero General Public License or a
separate commercial Swiss Ephemeris Professional license. A
developer who wants to build or redistribute a Vedic astrology tool under a different license,
or who wants to independently audit where a returned planetary position actually comes from,
is left with two options: adopt Swiss Ephemeris's AGPL obligations (or pay for its commercial
license), or hand-verify numerical tables against published references without any systematic
provenance trail. Neither option gives a permissively-licensable, provenance-auditable
alternative. Vedaksha was built to close that gap.

By ``clean-room'' we mean specifically this: every algorithm and constant in Vedaksha is derived
from a primary source --- a peer-reviewed analytical theory such as VSOP87A for the major
planets or ELP/MPP02 for the Moon \citep{chapront2003elpmpp02}, or a numerical kernel
distributed by a national agency, such as the JPL DE440s SPK ephemeris --- independently
validated against a JPL DE441 ephemeris served by the NASA/JPL Horizons web service, which is
never itself ingested as a local kernel --- and the
algorithms and constants shipped in the current release contain no source code, header file, or
documentation belonging to an existing proprietary or copyleft-licensed ephemeris implementation.
The project's history is more instructive than a blanket denial: three early ayanamsha
(sidereal-zodiac offset) constants were originally identified by reading a published header file
of the AGPL-licensed Swiss Ephemeris library to find a numeric value --- no source code was
copied, and the numeric constants themselves are uncopyrightable facts --- and all three were
subsequently re-derived from independently cited primary sources in a 2026-08 clean-room pass
that first removed the original values from the working tree, so the re-derivation could not see
its own target. This claim is deliberately made verifiable
rather than asserted on trust: \texttt{DATA\_PROVENANCE.md}, committed at the root of the
repository, lists every external data source the project ingests, together with its primary
URL, license status, fetch date, and a content hash where one applies; and the \texttt{docs/audit/}
directory carries dated, per-derivation clean-room records --- including a firewall-enforcement
transcript for the 2026-08-17 re-derivation of the ayanamsha (sidereal-zodiac offset) constants
--- that document what was and was not consulted at the time each derivation was made public.
These records are public alongside the code, not retrospective summaries, so a downstream user
can check the provenance claim against the same evidence this paper cites.

The remainder of this paper is organized as follows: Section~2 describes the engine's
architecture and the ephemeris and jyotish computation methods it implements; Section~3 reports
independent numerical validation against a public ephemeris oracle; Section~4 describes the
software's distribution channels, interfaces, and licensing; and Section~5 concludes.

\section{Methods}
PLACEHOLDER -- replaced in Task 4.

\section{Validation}
PLACEHOLDER -- replaced in Task 5.

\section{Software Architecture and Availability}
PLACEHOLDER -- replaced in Task 6.

\section{Conclusion}
PLACEHOLDER -- replaced in Task 6.

\bibliographystyle{plainnat}
\bibliography{references}

\end{document}
